\clearpage
\VerdierPhysicalPageAnchor{0122}%
\begingroup
\small
\setlength{\emergencystretch}{1.5em}%
\setlength{\abovedisplayskip}{4pt plus 2pt minus 2pt}%
\setlength{\belowdisplayskip}{4pt plus 2pt minus 2pt}%

\VerdierSemanticAnchor{ch2-1-2-13-countability-hypothesis}
\textit{$\mathcal D$ において可算積または可算直和が表現可能ならば、
圏 $\mathcal D$ は分解可能である。}

\VerdierSemanticAnchor{ch2-1-2-13-proof-morphism}
実際、$fv=0$ を満たす射 $v\colon W\longrightarrow X$ をとる。
このとき、次の射の射がある。

\VerdierSemanticAnchor{ch2-1-2-13-morphism-square}
\[
  \begin{tikzcd}[column sep=4em,row sep=2.5em]
    X \arrow[r,"f"] & Y \\
    W \arrow[u,"v"] \arrow[r] & 0 \arrow[u]
  \end{tikzcd}
  \,.
\]

\VerdierSemanticAnchor{ch2-1-2-13-1}
したがって関手 $\operatorname{Tr}$ を適用すると、次の可換図式を得る。
\begin{equation}
  \resizebox{0.82\linewidth}{!}{%
    $\begin{tikzcd}[ampersand replacement=\&,column sep=2.05em,row sep=3em]
      X \arrow[r,"f"]
        \& Y \arrow[r,"{a(f)}"]
        \& \tau(f) \arrow[r,"{b(f)}"]
        \& X[1] \\
      W \arrow[u,"v"] \arrow[r]
        \& 0 \arrow[u] \arrow[r]
        \& \tau(W\longrightarrow 0) \arrow[u,"{\operatorname{Tr}(v,0)}"]
          \arrow[r,"{b(W\longrightarrow 0)}"']
        \& W[1] \arrow[u,"{v[1]}"']
    \end{tikzcd}$%
  }\,.
  \tag{1.2.13.1}
\end{equation}

\VerdierSemanticAnchor{ch2-1-2-13-bottom-triangle}
まず、図式 (1.2.13.1) の下段は完全三角であることに注意する。
したがって射 $b(W\longrightarrow 0)$ は同型である (1.2.1)。そこで、

\VerdierSemanticAnchor{ch2-1-2-13-alpha-definition}
\[
  \alpha(v)=\operatorname{Tr}(v,0)[-1]
  b(W\longrightarrow 0)^{-1}[-1]\,,
\]
とおけば、図式

\VerdierSemanticAnchor{ch2-1-2-13-alpha-diagram}
\[
  \begin{tikzcd}[column sep=3.5em,row sep=2.4em]
    \tau(f)[-1] \arrow[r,"{b(f)[-1]}"]
      & X \arrow[r,"f"]
      & Y \\
      & W \arrow[ul,"{\alpha(v)}"'] \arrow[u,"v"] \arrow[ur,"0"']
  \end{tikzcd}
\]
は可換である。このようにして次の切断が定まることは直ちに確かめられる。
\par

\VerdierSemanticAnchor{ch2-1-2-13-alpha-section}
\[
  \alpha\colon \operatorname{Ker}(f)\longrightarrow \tau(f)[-1]
\]
\noindent これは、$\tau(f)[-1]$ が表現する関手から射 $f$ の核関手への標準射
$\tau(f)[-1]\longrightarrow\operatorname{Ker}(f)$ の切断である。

\VerdierSemanticAnchor{ch2-1-2-13-projector-kernel}
したがって $f$ の核関手は、$\tau(f)[-1]$ のある射影子の核関手と同型である。
ゆえに、この核関手は表現可能である (1.2.9)。

\VerdierSemanticAnchor{ch2-1-2-13-proof-complete}
(1.2.9) により射 $f$ は分解可能である。これで証明が完了した。
\endgroup
